\scanpage{132}
% archive sheet number: 131
\uline{Corollaire}\quad Fonctorialité \tikz[baseline=(a.base)]{\node[draw,ellipse,inner sep=2pt] (a) {en $X_K$};} de l'ensemble $\operatorname{\text{Ép}}\deleted{\unsure{\mathrm{ais}}}(\bar X)$ des points \deleted{\unsure{de}} épais de $\bar X$ [munis de leurs corps résiduels, ----], pour des \deleted{\unsure{morphismes}}\inserted{applic. rationnelles} $X_K\unsure{\mathord{-}|\longrightarrow}Y_K$ \uline{dominantes gén.\textsuperscript{t} séparables de} $K$-\uline{schémas} \uline{propres}

\uline{Proposition} 2\quad Soit $S$ une partie constructible de $\bar X$, $X_K$-simple, [tel que $\unsure{l(N,x)|l(V,x)|(V,x)}$ y soit borné], et tel que $\bar\varphi(S)$ existe [p.ex. $Y$ propre sur $V$] et soit $Y_K$-simple. Alors $\bar\varphi(S)$ est constructible. De plus, il y a une partition \inserted{finie} $S=\coprod_i S_i$, les $S_i$ des sous-schémas irréductibles constructibles de $X$, \unsure{telles|telle} que $\varphi$ soit morphique d'indice fini sur chaque $S_i$, [et $\varphi(S_i)$ est alors un ens constructible irréductible \ldots]

\uline{Corollaire}\quad Fonctorialité de $X_K$ variant comme dessus, et de plus simple sur $K$, de l'ens $\operatorname{Cons}(\bar X)$ des parties constructibles de $\bar X$.
